%% =====================================================================
%%  T-19-01  The Uninvited Debt
%%  -------------------------------------------------------------------
%%  One idea per file. Core is mandatory. Slots in canonical order.
%%  Max 700 words. No layout commands. No filler. New source material
%%  for this entry goes in sources/<BOOK>/T-19-01.tex -- never by
%%  rewriting this file (docs/RULES.md R0.1).
%% =====================================================================
%% @kind: topic
%% @id: T-19-01
%% @title: The Uninvited Debt
%% @subtitle: An obligation does not require your consent
%% @chapter: c19
%% @order: 10
%% @status: stable
%% @sources: B4
%% @updated: 2026-09-19
%% @allow-rewrite: slot migration to the seven-question format (docs/RULES.md section 2)

\topic{T-19-01}{The Uninvited Debt}{An obligation does not require your consent}


\begin{core}
A benefit you never asked for, from someone you never chose, creates the same pressure to repay as one you negotiated --- and usually more, because declining it reads as ingratitude rather than as a decision. Consent is not part of the mechanism.
\end{core}

\begin{mechanism}
Reciprocity is enforced socially rather than contractually. The person who takes a benefit and returns nothing is marked, and the marking is public. That makes the rule hard to opt out of, because the opt-out is itself visible and is read as a verdict on character.

So the operative fact is not the value of the benefit but its occurrence. An unrequested favour installs an account whose balance was never agreed, whose repayment date is chosen by the creditor, and whose currency the creditor will nominate later. Critically, the repayment need have nothing to do with the original benefit: the account is not a trade, it is a standing obligation, and it can be drawn on for any request the holder likes to make. Choosing the form and timing of the benefit is what gives the holder control over the draw.
A clipboard, a flower pressed into a hand, an audit offered free of charge: the benefit is chosen so that declining it in public costs more than accepting it.
\end{mechanism}


\begin{application}
  \item Give first, on your own terms, and keep accounts small and closed; an open balance is an asset to whoever holds it.
  \item Choose a benefit that is genuinely useful and unmistakably voluntary, so that its reality cannot later be disputed.
  \item Decide the request you will make before the benefit is given, and make it in a different currency from the benefit.
  \item Never name the debt. Naming converts an asset into a grievance and releases the other party from it.
\end{application}

\begin{feedback}
  \item A benefit arrives that you would have declined had you been asked.
  \item The giver chose its form, its size and its timing.
  \item The request that follows is unrelated to the benefit and larger than it.
  \item Refusing produces a feeling about yourself rather than a calculation about the transaction.
  \item The benefit is mentioned --- by the giver or by a third party --- shortly before the request.
\end{feedback}

\begin{failure}
Refusing every unsolicited benefit produces a person who cannot accept help and reads all generosity as an opening move, which is its own captivity and costs real relationships. The rule worth holding is narrow: scrutinise where the benefit was unrequested, large relative to the relationship, or followed quickly by an unrelated ask. Elsewhere let the account stand --- ordinary reciprocity is what makes cooperation affordable in the first place.
\end{failure}

\begin{countermeasures}
  \item Refuse early and politely, in the same breath as the offer. Once acceptance has been observed, the only remaining question is the size of the repayment, and the other party sets it.
  \item Where refusal is impossible, repay at once, specifically, and in a form you choose. A closed account cannot be drawn on.
  \item Split the two questions --- \emph{was I helped?} and \emph{does this request follow from it?} --- and answer the second on its own.
  \item Price the benefit at what it cost the giver, not at what it is worth to you.
  \item Make the refusal short and unexplained. A reason offered is a reason to be argued with.
\end{countermeasures}

\sources{B4}
\seesources
\seealso{T-04-02, T-19-02, T-24-01, T-16-01}
